\documentclass[a4paper,12pt,oneside,theme=amber]{maki-notes}

\renewcommand{\LectureNotesTitle}{TikZ Style Gallery}
\renewcommand{\AuthorInfo}{Test Runner}

\newcommand{\ThemeCard}[1]{%
  \begingroup
  \csname makiapplytheme#1\endcsname
  \begin{tikzpicture}[x=1cm,y=1cm,baseline=(base.center)]
    \node[inner sep=0pt] (base) at (0,0) {};
    \path[draw=MainText!18, fill=white, rounded corners=2.5pt] (0,0) rectangle (4.25,1.2);
    \fill[LogicColor] (0.22,0.78) rectangle +(0.74,0.18);
    \fill[DefColor]   (1.08,0.78) rectangle +(0.74,0.18);
    \fill[ExColor]    (1.94,0.78) rectangle +(0.74,0.18);
    \fill[NoteColor]  (2.80,0.78) rectangle +(0.74,0.18);
    \fill[GeometryColor] (3.46,0.78) rectangle +(0.57,0.18);
    \node[anchor=west, font=\ttfamily\footnotesize, text=MainText] at (0.22,0.38) {#1};
    \node[anchor=east, font=\footnotesize, text=MainText!78] at (4.03,0.38) {\MakiResolvedTheme};
  \end{tikzpicture}%
  \endgroup
}

\begin{document}
\mainmatter

\chapter{TikZ Style Gallery}

\section{Theme Presets}

这一页同时验证主题色位与 TikZ 预设族是否能一起工作. 文档类本身使用 `theme=amber`, 下方色卡则额外切换并展示所有当前支持的主题, 这样一份测试文档就能覆盖新增主题入口和色位联动.

\begin{center}
\ThemeCard{default}\hspace{0.6em}
\ThemeCard{ocean}\hspace{0.6em}
\ThemeCard{forest}

\medskip

\ThemeCard{graphite}\hspace{0.6em}
\ThemeCard{amber}\hspace{0.6em}
\ThemeCard{berry}

\medskip

\ThemeCard{sandstone}
\end{center}

\section{Deep Learning}

这一页用于验证神经网络相关样式是否能在讲义中直接复用. 布局思路参考了 tikz.net 的 \texttt{neural\_networks} 页面: 节点样式按输入层、隐藏层、输出层分色, 连线、跳连与张量块则拆成语义化样式, 方便在不同网络结构中复用.

\begin{center}
\begin{tikzpicture}[x=1.45cm, y=1.0cm]
\foreach \name/\y/\label in {I1/1.6/$x_1$, I2/0.55/$x_2$, I3/-0.55/$x_3$, I4/-1.6/$x_4$}
  \node[nn input] (\name) at (0,\y) {\label};

\foreach \name/\y in {H11/2.1, H12/0.7, H13/-0.7, H14/-2.1}
  \node[nn hidden] (\name) at (1.9,\y) {};
\foreach \name/\y in {H21/1.6, H22/0.4, H23/-0.8, H24/-2.0}
  \node[nn hidden] (\name) at (3.8,\y) {};

\node[nn output] (O1) at (5.8,0.9) {$\hat y_1$};
\node[nn output] (O2) at (5.8,-0.9) {$\hat y_2$};

\foreach \i in {1,2,3,4}{
  \foreach \j in {1,2,3,4}{
    \draw[nn connect] (I\i) -- (H1\j);
    \draw[nn connect] (H1\i) -- (H2\j);
  }
}
\foreach \j in {1,2,3,4}{
  \draw[nn connect strong] (H2\j) -- (O1);
  \draw[nn connect strong] (H2\j) -- (O2);
}

\node[nn tensor] (E) at (-1.5,0) {embed};
\draw[nn attention] (E.east) -- (I2.west);
\draw[nn attention] (E.east) -- (I3.west);

\node[nn conv] (C1) at (2.85,3.25) {Conv\\$3\times 3$};
\draw[nn skip] (H12.north) to[out=90, in=180] (C1.west);
\draw[nn skip] (C1.east) to[out=0, in=90] (H22.north);

\node[nn sum] (S) at (4.8,0) {$+$};
\draw[nn connect strong] (H22.east) -- (S.west);
\draw[nn skip] (I2.east) to[out=0, in=180] (S.west);
\draw[nn connect strong] (S.east) -- ($(O1)!0.5!(O2)$);

\node[nn block, fit=(H11)(H14), label={[maki label]above:Feature extractor}] {};
\node[nn block, fit=(H21)(H24), label={[maki label]above:Task head}] {};
\node[maki annotation] at ($(O1)!0.5!(O2)+(1.6,0)$) {classification\\scores};
\end{tikzpicture}
\captionof{figure}{多层感知机与跳连的样式测试}
\end{center}

\section{Mathematics and Geometry}

这一部分参考了 tikz.net 上 \texttt{osculating-circle} 与 \texttt{vector-field} 一类图的组织方式: 曲线、采样点、切线、法线、向量场箭头分别归为不同语义样式, 这样写分析几何、微分几何或数学习题讲义时不用每次手写绘图参数.

\begin{center}
\begin{tikzpicture}[scale=0.95]
\draw[math axis] (-0.5,0) -- (6.8,0) node[right] {$x$};
\draw[math axis] (0,-1.0) -- (0,4.8) node[above] {$y$};
\draw[maki guide] (0,0) grid[xstep=1, ystep=1] (6,4);

\fill[math region] plot[smooth, domain=0.6:5.0, samples=100]
  (\x,{0.15*(\x-1.1)*(\x-3.7)*(\x-5.4)+2.7}) -- (5.0,0) -- (0.6,0) -- cycle;
\draw[math curve] plot[smooth, domain=0.6:5.6, samples=120]
  (\x,{0.15*(\x-1.1)*(\x-3.7)*(\x-5.4)+2.7});

\coordinate (P) at (3.1,{0.15*(3.1-1.1)*(3.1-3.7)*(3.1-5.4)+2.7});
\node[math sample point, label={[maki label]above right:$P$}] at (P) {};
\draw[math tangent] ($(P)+(-1.4,-0.6)$) -- ($(P)+(1.6,0.68)$) node[pos=0.93, above] {$T_P\gamma$};
\draw[math normal] ($(P)+(-0.45,1.3)$) -- ($(P)+(0.45,-1.3)$) node[pos=0.1, right] {$N_P\gamma$};

\foreach \x/\y in {0.9/0.55,1.5/1.1,2.2/1.6,2.9/2.05,3.7/2.15,4.6/1.5,5.3/0.9}{
  \pgfmathsetmacro{\vx}{0.24*(0.6+\y)}
  \pgfmathsetmacro{\vy}{0.16*(2.6-\x)}
  \draw[math vector] (\x,\y) -- ++(\vx,\vy);
}

\node[maki annotation] at (5.35,3.85) {curve\\tangent\\normal\\vector field};
\end{tikzpicture}
\captionof{figure}{分析与微分几何风格的样式测试}
\end{center}

\section{Topology and Geometry}

最后一组样式借用了 tikz.net \texttt{plane-to-torus} 的基本域识别箭头写法. 这里重点不是完全复刻原图, 而是把“边识别”“曲面边界”“基本环路”“地质线/切割线”这些在拓扑与几何讲义中高频出现的视觉元素抽出来.

\begin{center}
\begin{tikzpicture}[x=1cm, y=1cm]
\begin{scope}
  \draw[topo patch] (0,0) rectangle (4.2,4.2);
  \foreach \k in {0.7,1.4,2.1,2.8,3.5}{
    \draw[topo mesh] (\k,0) -- (\k,4.2);
    \draw[topo mesh] (0,\k) -- (4.2,\k);
  }
  \draw[topo identify] (0,0) -- (4.2,0);
  \draw[topo identify] (0,4.2) -- (4.2,4.2);
  \draw[topo identify reversed] (0,0) -- (0,4.2);
  \draw[topo identify reversed] (4.2,0) -- (4.2,4.2);
  \draw[topo geodesic] (0.5,0.8) .. controls (1.2,2.2) and (2.7,2.9) .. (3.8,3.5);
  \draw[topo cut] (1.0,3.7) .. controls (1.7,2.8) and (2.7,1.4) .. (3.4,0.5);
  \node[maki annotation] at (2.1,2.1) {fundamental\\domain};
  \node[maki label] at (2.1,-0.45) {$a$};
  \node[maki label] at (2.1,4.65) {$a$};
  \node[maki label] at (-0.35,2.1) {$b^{-1}$};
  \node[maki label] at (4.55,2.1) {$b^{-1}$};
\end{scope}

\begin{scope}[xshift=7.4cm]
  \draw[topo boundary, fill=TopologyBg] (0,0) ellipse (2.15 and 1.35);
  \draw[topo boundary, fill=white] (0,0) ellipse (0.78 and 0.42);
  \draw[topo loop] (0,0.72) ellipse (1.75 and 0.42);
  \draw[topo geodesic] (-1.2,-0.05) .. controls (-0.3,1.15) and (0.95,1.0) .. (1.35,0.08);
  \draw[topo cut] (0,-1.35) -- (0,1.35);
  \node[maki annotation] at (0,-2.0) {meridian / longitude\\style sketch};
\end{scope}
\end{tikzpicture}
\captionof{figure}{基本域、边识别与环面的拓扑样式测试}
\end{center}

\section{Flowcharts and Timelines}

这一组新增样式参考了 tikz.net 上 flowchart 与 timeline 一类图的组织方法: 节点形状、流程箭头、里程碑和时间窗口各自独立命名, 以后画算法流程、课程安排或项目进度时就不需要重新拼装.

\begin{center}
\begin{tikzpicture}[x=1cm, y=1cm]
\node[flow terminal] (start) at (0,2.1) {开始};
\node[flow process] (readme) at (2.8,2.1) {读取题面\\与仓库约束};
\node[flow decision] (branch) at (5.9,2.1) {需要\\新图形族?};
\node[flow process] (style) at (8.8,2.1) {整理语义化\\TikZ 样式};
\node[flow io] (docs) at (8.8,0.2) {更新 README\\与测试文档};
\node[flow terminal] (done) at (11.7,2.1) {完成};
\node[flow note, anchor=west] at (9.8,0.2) {预设类目: flow / timeline / probability / graph};

\draw[flow arrow] (start) -- (readme);
\draw[flow arrow] (readme) -- (branch);
\draw[flow arrow emph] (branch) -- node[above, maki label] {是} (style);
\draw[flow arrow] (style) -- (done);
\draw[flow arrow feedback] (branch) -- ++(0,-1.55) -| node[pos=0.25, left, maki label] {补文档} (docs.west);
\draw[flow arrow] (docs.north) |- (style.south);

\begin{scope}[yshift=-3.55cm]
  \draw[timeline axis] (0,0) -- (11.7,0);
  \foreach \x/\year in {0.9/2024,3.1/2025,5.7/2026,8.2/2027,10.8/2028}{
    \draw[timeline tick] (\x,-0.15) -- (\x,0.15);
    \node[maki label, below] at (\x,-0.18) {\year};
  }
  \draw[timeline span] (1.15,0.48) -- (5.45,0.48);
  \node[timeline milestone] (m1) at (2.15,1.1) {模板初版};
  \node[timeline event] (m2) at (5.7,0) {};
  \node[timeline milestone] (m3) at (8.2,1.1) {主题扩展};
  \node[timeline event, fill=LogicBg, draw=LogicColor!80!black] (m4) at (10.8,0) {};
  \draw[timeline connector] (m1.south) -- (2.15,0.18);
  \draw[timeline connector] (m3.south) -- (8.2,0.18);
  \node[maki label, above=2pt of m2] {wrapfig};
  \node[maki label, above=2pt of m4] {release};
\end{scope}
\end{tikzpicture}
\captionof{figure}{流程图与时间轴预设样式测试}
\end{center}

\section{Probability and Graph Structures}

最后一组新增样式面向概率论与图论讲义. 参考了 tikz.net 的 Venn/集合图和 graph-isomorphism 一类节点边图, 把样本空间、事件区域、条件区域、顶点和特殊边都抽成可复用样式.

\begin{center}
\begin{tikzpicture}[x=1cm, y=1cm]
\begin{scope}
  \draw[prob universe] (0,0) rectangle (5.5,4.2);
  \node[maki label, above left] at (0,4.2) {$\Omega$};
  \filldraw[prob event] (1.45,2.15) ellipse (1.1 and 1.5);
  \filldraw[prob event alt] (3.9,2.75) ellipse (0.82 and 0.72);
  \filldraw[prob event alt] (4.0,1.15) ellipse (0.95 and 0.72);
  \draw[prob condition] (2.75,0) -- (2.75,4.2);
  \foreach \x/\y in {0.95/2.95,1.2/2.3,1.55/1.8,1.8/2.55,1.45/1.2,3.6/2.95,4.15/2.6,3.75/1.15,4.25/0.95,4.55/1.3}{
    \node[prob outcome] at (\x,\y) {};
  }
  \node[prob label] at (1.45,3.7) {$A$};
  \node[prob label] at (4.15,3.45) {$B$};
  \node[prob label] at (4.15,1.85) {$C$};
  \node[prob label, anchor=west] at (2.92,3.95) {$A^c$};
\end{scope}

\begin{scope}[xshift=7.2cm]
  \node[graph node accent] (v1) at (0,1.9) {};
  \node[graph node] (v2) at (1.8,3.0) {};
  \node[graph node] (v3) at (3.8,2.6) {};
  \node[graph node muted] (v4) at (3.9,0.9) {};
  \node[graph node] (v5) at (1.7,0.1) {};
  \node[graph node] (v6) at (0.1,0.7) {};
  \foreach \u/\v/\label in {
    v1/v2/2, v2/v3/1, v3/v4/4, v4/v5/3, v5/v6/2, v6/v1/1, v1/v5/5, v2/v5/2
  }{
    \draw[graph edge] (\u) -- (\v) node[midway, graph label, fill=white] {\label};
  }
  \draw[graph edge strong] (v1) -- (v3);
  \draw[graph edge cut] (v2) -- (v4);
  \draw[graph walk] (v6) to[bend left=12] (v3);
  \foreach \name/\text in {v1/s, v2/a, v3/b, v4/t, v5/c, v6/d}{
    \node[maki label] at (\name) {\text};
  }
  \node[flow note, anchor=west] at (4.55,2.9) {accent edge = shortest path\\dashed edge = cut candidate};
\end{scope}
\end{tikzpicture}
\captionof{figure}{概率/集合图与图论节点边样式测试}
\end{center}

\end{document}
